%%%%%%%%%%%%%%%%%%%%%%%%%%%%%%%%%%%%%%%%%%%%%%%%
% chapter4.tex                                 %
% Contains formatting and content of chapter 4 %
%%%%%%%%%%%%%%%%%%%%%%%%%%%%%%%%%%%%%%%%%%%%%%%%
\chapter{Conclusion}
\newpage

\section{Conclusion}

Using the Binomial-Beta model with the Slippi data and varying priors, we determined that the Marth-Fox matchup is close enough to be considered even. This contrasts Leffen's claim that the Marth-Fox matchup is heavily favored for Marth. With the various logistic regression models, using a subset of the Slippi data, we determined that Yoshi's Story is definitively a strong stage for Marth against Fox. While not definitive, Final Destination appears to be in Marth's favor, while Dream Land and Pokemon Stadium appear to lean in Fox's favor. Battlefield and Fountain of Dreams appear to be relatively even. Final Destination is often regarded as Marth's best stage in the matchup, so it was interesting to see the average Marth player perform better on Yoshi's Story, relatively speaking. Finally, we adapted the logistic regression model to factor in the relative difference of player Slippi ratings. Although this data was not in our dataset, future Slippi match data containing player ratings could use such data in this model for analysis.

\indent There are nearly infinite variables when it comes to fighting games. A matchup may be awful on paper, but a player may practice it to such a degree that they make seem as though their character wins a matchup. Additionally, in competitive Melee, no two games are truly independent. Players use information about their opponents playstyles, habits, and skills to make strategic, on-the-fly adaptations. The free-flowing nature of a game like Melee results in infinite possibilities, and it is a main contributer to why its competitive scene has not only survived, but thrived over the last 25 years.

\indent For the scope of this project, I opted to make assumptions about match independence and select a few major matchup factors to simplify the creation of a logistic regression model. There are several steps to improve upon the existing model. First, the model could utilize actual tournament data, where we could identify win probabilities given the result of the current match count in a best-of-five set. Additionally, data could be polled for specific Marth players to see how they do against specific Fox players or Fox players in general. A similar approach could be used for other characters in the cast. We could also run the model using data from only top players to determine if matchups change at different skill levels, or we could introduce an interaction effect to see if a combination of skill difference and stage selection work together to affect a matchup. Analysis of data like this could improve the viewer experience by showing live probabilities of a player winning given the current state of a set, or can be used to predict winners of major tournaments. Like Melee itself, there are nearly infinite possible ways to leverage Baysian methods to uncover something new about a beloved game.

